% ============================================================
% preamble.tex —— 宏包与全局样式定义
% 团队成员一般不需要修改本文件；如需新增宏包请在此处添加，
% 并在 README.md 的“依赖说明”中登记，保证全组编译环境一致。
% ============================================================

% ---------- 页面与版式 ----------
\usepackage[a4paper,top=2.5cm,bottom=2.5cm,left=2.6cm,right=2.6cm]{geometry}
\usepackage{fancyhdr}
\pagestyle{fancy}
\fancyhf{}
\fancyhead[L]{\small 2026 Agent大赛《驾驶安全》算法赛道}
\fancyhead[R]{\small 任务一：事故预测算法说明文档}
\fancyfoot[C]{\small \thepage}
\renewcommand{\headrulewidth}{0.4pt}

% ---------- 数学与符号 ----------
\usepackage{amsmath,amssymb}

% ---------- 表格 ----------
\usepackage{booktabs}      % 三线表
\usepackage{longtable}     % 跨页长表
\usepackage{tabularx}      % 自动换行列
\usepackage{multirow}
\usepackage{array}

% ---------- 图形 ----------
\usepackage{graphicx}
\graphicspath{{figures/}}  % 图片统一放 figures/ 目录，引用时直接写文件名
\usepackage{float}

% ---------- 颜色与提示框 ----------
\usepackage{xcolor}
\usepackage{tcolorbox}
\tcbuselibrary{breakable}  % 提示框可跨页断开

% ---------- 列表 ----------
\usepackage{enumitem}
\setlist{nosep}

% ---------- 图注表注 ----------
\usepackage[font=small,labelfont=bf]{caption}

% ---------- 代码 ----------
\usepackage{listings}
\lstset{
  basicstyle=\small\ttfamily,
  breaklines=true,
  columns=flexible,
  numbers=none,
  frame=single,
  framerule=0.3pt,
  rulecolor=\color{gray!60},
  backgroundcolor=\color{gray!5}
}

% ---------- 超链接 ----------
\usepackage[colorlinks,linkcolor=blue!60!black,citecolor=blue!60!black,urlcolor=blue!60!black]{hyperref}

% ============================================================
% 自定义命令
% ============================================================

% 待填写提示框（环境用法）：
%   \begin{TODO}
%   说明文字……
%   \end{TODO}
% 各章节中所有 TODO 框在最终提交前必须全部删除。
\newtcolorbox{TODO}{colback=yellow!10,colframe=orange!80!black,
  boxrule=0.6pt,arc=1.5mm,left=2mm,right=2mm,top=1mm,bottom=1mm,
  breakable,fontupper=\small,
  before upper={\textbf{【待填写】}\ }}

% 重点结论框（环境用法）：\begin{hlbox}...\end{hlbox}
\newtcolorbox{hlbox}{colback=blue!4,colframe=blue!55!black,
  boxrule=0.6pt,arc=1.5mm,left=2mm,right=2mm,top=1.5mm,bottom=1.5mm,breakable}

% 示例数据框（等宽字体展示原始记录，保留源码换行、左对齐）
\newtcolorbox{databox}[1][]{colback=gray!5,colframe=gray!55,
  boxrule=0.5pt,arc=1mm,left=2mm,right=2mm,top=1.5mm,bottom=1.5mm,
  fontupper=\small\ttfamily\obeylines\raggedright,#1}

% 常用表头样式
\newcolumntype{Y}{>{\centering\arraybackslash}X}
\newcolumntype{L}{>{\raggedright\arraybackslash}X}
